\subsection{Extension of homotopies}

PROPOSITION 1. --- Let X' be a normal topological space, let X be a subspace of X', let A be a closed subspace of X' contained in X, and let U be a neighborhood of A in X.

Denote by \(i_{A}\) the canonical injection of A into X, \(i_{U}\) the canonical injection from U into X; denote by \(j_{A}\) and \(j_{U}\) the corresponding injections from \(\mathrm{Cyl}(i_{\mathrm{A}})\) and \(\mathrm{Cyl}(i_{\mathrm{U}})\) into \(X \times I\).

There exists a continuous map \(r \colon \mathrm{X} \times \mathbf{I} \to \operatorname{Cyl}(i_{\mathrm{U}})\) such that \(j_{\mathrm{U}} \circ r(x) = x\) for every point \(x \in j_{\mathrm{A}}(\operatorname{Cyl}(i_{\mathrm{A}}))\) .

Let U' be an open subset of X' such that A ⊂ X ∩ U' ⊂ U. By definition of a normal space (TG, IX, p. 41), there exists an open neighborhood V' of A in X' such that A ⊂ V' ⊂ \(\overline{V'}\) ⊂ U' and a continuous function φ': X' → I that is equal to 1 at every point of A and to 0 at every point of CV'. Set φ = φ' \textbar{} X and V = X ∩ V'. Also denote α: U × I → Cyl(i\_U) and β: X → Cyl(i\_U) the canonical maps (III, p. 238).

Let \(r\) be the map from \(\mathrm{X} \times \mathbf{I}\) into \(\operatorname{Cyl}(i_{\mathrm{U}})\) given by

\[
r (x, t) = \left\{ \begin{array}{l l} \alpha (x, 1 - (1 - t) \varphi (x)) & \text { for } (x, t) \in \mathrm{U} \times \mathbf {I}, \\ \beta (x) & \text { otherwise }. \end{array} \right.
\]

The map \(r\) is continuous on \(\mathrm{U} \times \mathbf{I}\). For every point \((x, t) \in \mathrm{U} \times \mathbf{I}\) such that \(x \notin \mathrm{V}\), we have \(\varphi(x) = 0\), whence \(r(x, t) = \alpha(x, 1) = \beta(x)\). Consequently, the maps \(r\) and \(\beta \circ \mathrm{pr}_1\) coincide on \((\mathrm{X} - \mathrm{V}) \times \mathbf{I}\), which implies that \(r\) is continuous on this subspace. Since U and the interior of \(\mathrm{X} - \mathrm{V}\) cover X, the map \(r\) is continuous.

Let \(y\) be a point of \(\mathrm{Cyl}(i_{\mathrm{A}})\); set \((x,t) = j_{\mathrm{A}}(y)\). If \(x \in \mathrm{A}\), \(\varphi(x) = 1\) and \(r(x,t) = \alpha(x,t)\), whence \(j_{\mathrm{U}}(r(x,t)) = (x,t)\). Otherwise, \(t = 1\) and one has \(r(x,t) = \alpha(x,1)\) if \(x \in \mathrm{U}\) and \(r(x,t) = \beta(x)\) otherwise; consequently, \(j_{\mathrm{U}}(r(x,t)) = (x,t)\) in this case. This implies that \(j_{\mathrm{U}} \circ r\) maps every point of \(j_{\mathrm{A}}(\mathrm{Cyl}(i_{\mathrm{A}}))\) to itself, whence the proposition.

COROLLARY 1. --- Let X be a normal topological space, let f be a continuous map from X into a topological space Z. Let A be a closed subspace of X, let U be a neighborhood of A in X, and let σ: U × I → Z be a homotopy whose terminal map is the map f \textbar{} U. There exists a homotopy \(\widetilde{\sigma}\colon\mathbf{X}\times\mathbf{I}\to\mathbf{Z}\) whose terminal map is \(f\) and which coincides with \(\sigma\) on \(\mathbf{A}\times\mathbf{I}\).

Let us set \(X' = X\) and let \(r: X \times I \to \text{Cyl}(i_{\text{U}})\) be the continuous map given by Prop. 1. With the notation introduced in the proof of this proposition, there exists a unique continuous map \(F: \text{Cyl}(i_{\text{U}}) \to Z\) such that \(\text{F}(\alpha(x,t)) = \sigma(x,t)\) for \((x,t) \in \text{U} \times \text{I}\) and \(\text{F}(\beta(x)) = f(x)\) for \(x \in X\) (III, p. 238, prop. 4). The map \(F \circ r: X \times I \to Z\) is a homotopy; its terminal map sends a point \(x \in X\) to \(\text{F}(r(x,1)) = \text{F}(\beta(x)) = f(x)\). If \((x,t) \in A \times I\), \(\text{F}(r(x,t)) = \text{F}(\alpha(x,t)) = \sigma(x,t)\); hence this homotopy coincides with \(\sigma\) on \(A \times I\).

COROLLARY 2. --- Let X be a normal topological space, let A be a closed subspace of X, and let U be a neighborhood of A in X. Let f and f' be continuous homotopic maps from U into a topological space Z. If f has a continuous extension to X, then so does the map f' \textbar{} A.

Let F be a continuous mapping from X into Z such that \(F \textbar{} U = f\). Let us choose a homotopy \(\sigma\colon A \times I \to Z\) joining f'	extbar{} A to f	extbar{} A. By Cor. 1, there exists a homotopy \(\widetilde{\sigma}\colon X \times I \to Z\) whose terminal map is F and which extends \(\sigma\). The initial map of \(\widetilde{\sigma}\) is then a continuous map from X to Z that coincides with \(f'\) on A.

COROLLARY 3. — Let X be a normal topological space, let A be a closed subspace of X, and let U be a neighborhood of A in X. Let Z be a topological space homotopy equivalent to a point, and let g: U → Z be a continuous map. There exists a continuous map \(\widetilde{g}\) : X → Z which coincides with g on A.

Since Z is homotopic to a point, the map g is homotopic to a constant map f. The map f extends continuously to X; the assertion therefore follows from Cor. 2.

COROLLARY 4. --- Let X be a paracompact topological space, let A be a closed subspace of X. Let n be an integer ≥ 0 and let V be an open subset of R \(^{n}\). Let f: X → V be a continuous map and let σ: A × I → V be a homotopy with terminal map f \textbar{} A. There exists a homotopy \(\widetilde{\sigma}\): X × I → V with terminal map f that extends σ.

Denote by \(i_{\mathrm{A}}\) the canonical injection of A into X, \(\alpha_{\mathrm{A}}: \mathrm{A} \times \mathbf{I} \to \operatorname{Cyl}(i_{\mathrm{A}})\) and \(\beta_{\mathrm{A}}: \mathrm{X} \to \operatorname{Cyl}(i_{\mathrm{A}})\) the canonical maps, as well as \(j_{\mathrm{A}}: \operatorname{Cyl}(i_{\mathrm{A}}) \to \mathrm{X} \times \mathbf{I}\) the canonical injection. Since A is closed in X, \(j_{\mathrm{A}}\) defines a homeomorphism from \(\mathrm{Cyl}(i_{\mathrm{A}})\) onto the closed subspace \((\mathrm{A}\times\mathbf{I})\cup(\mathrm{X}\times\{1\})\) of \(X\times I\). Let \(\widetilde{\sigma}_{0}\) be the unique map from \(\mathrm{Cyl}(i_{\mathrm{A}})\) into V such that \(\widetilde{\sigma}_{0}\circ\alpha_{A}=\sigma\) and \(\widetilde{\sigma}_{0}\circ\beta_{A}=f\); it is continuous (III, p. 238, prop. 4).

Since X is paracompact and I is compact, the space \(X \times I\) is paracompact (TG, I, p. 70, prop. 17), hence normal (TG, IX, p. 49, prop. 4). There therefore exists a continuous map \(k: X \times I \to R^{n}\) such that \(k \circ j_{A} = \widetilde{\sigma}_{0}\) (TG, IX, p. 45, corollary).

The set U of points \(x \in X\) such that \(k(x, t) \in V\) for all \(t \in I\) is open in X (IV, p. 439, lemma 1). It is therefore a neighborhood of A. By Corollary 1, applied to the map f and to the homotopy \(k | U \times I\), there exists a homotopy \(\widetilde{\sigma} \colon X \times I \to V\) with initial map f which coincides with k, hence with \(\sigma\), on \(A \times I\).

Lemma 1. — Let Z be a topological space, K a compact topological space, and W an open subset of Z × K. The set U of points z ∈ Z such that \{z\} × K is contained in W is open in Z.

The first projection \(pr_{1}: Z \times K \to Z\) is proper (TG, I, p. 77, cor. 5), hence closed. Consequently, \(\mathbb{C}U = \mathrm{pr}_{1}(\mathbb{C}W)\) is closed in Z and U is open.

COROLLARY 5. --- Let X' be a normal topological space, let X be a subspace of X', let A be a closed subspace of X' contained in X, and let U be a neighborhood of A in X.

Let Y and Z be topological spaces; let \(f\colon X \times \{1\} \times Y \to X \times \{1\} \times Z\) be an \(X \times \{1\}\)-morphism and let \(g\colon U \times I \times Y \to U \times I \times Z\) be a \(U \times I\)-morphism that agrees with \(f\) on \(U \times \{1\} \times Y\).

There then exists an \(X \times I\)-morphism \(h: X \times I \times Y \to X \times I \times Z\) which coincides with f on \(X \times \{1\} \times Y\) and with g on \(A \times I \times Y\).

Moreover, if f and g are homeomorphisms, one can choose a homeomorphism h having the required properties.

Let U' be an open neighborhood of A in X' such that U'∩X ⊂ U. Since the space X' is normal, there exists an open neighborhood V' of A in X' such that A ⊂ V' ⊂ \(\overline{V'}\) ⊂ U' (TG, IX, p. 41). Replacing U if necessary by \(\overline{V'}\) ∩ X, one may thus suppose that U is closed in X. We then take up again the notations of prop. 1 and its proof; let \(r: X \times I \to \text{Cyl}(i_{\text{U}})\) be a continuous map such that \(j_{\text{U}} \circ r(x) = x\) for every point \(x \in j_{\text{A}}(\text{Cyl}(i_{\text{A}}))\).

Also put \(f' = pr_{3} \circ f\) and \(g' = pr_{3} \circ g\) . Since \(f'\) and \(g'\) coincide on \(U \times \{1\} \times Y\) , there exists a unique map

\[
\varphi \colon \operatorname{Cyl} (i _ {\mathrm{U}}) \times \mathrm{Y} \to \operatorname{Cyl} (i _ {\mathrm{U}}) \times \mathrm{Z}
\]

such that \(\varphi(\alpha(x,t),y) = (\alpha(x,t),g'(x,t,y))\) for \((x,t,y) \in \mathrm{U} \times \mathbf{I} \times \mathrm{Y}\) and \(\varphi(\beta(x),y) = (\beta(x),f'(x,1,y))\) for \((x,y) \in \mathrm{X} \times \mathrm{Y}\). Since U is closed in X, the canonical surjection \(\pi\) from \((\mathrm{U} \times \mathbf{I}) \cup \mathrm{X}\) onto \(\mathrm{Cyl}(i_{\mathrm{U}})\) is universally strict (III, p. 239, remark 2). Since the map \(\varphi \circ (\pi \times \mathrm{Id}_{\mathrm{Y}})\) is continuous, the map \(\varphi\) is continuous. It is therefore a morphism of \(\mathrm{Cyl}(i_{\mathrm{U}})\)-spaces, and even an isomorphism if \(f\) and \(g\) are homeomorphisms.

Now let \(h \colon \mathrm{X} \times \mathbf{I} \times \mathrm{Y} \to \mathrm{X} \times \mathbf{I} \times \mathrm{Z}\) be the map \(r^*(\varphi)\) deduced from \(\varphi\) by base change \(r\). It is given by \(h(x,t,y) = (x,t,\mathrm{pr}_2(\varphi(r(x,t),y)))\) for \((x,t,y) \in \mathrm{X} \times \mathbf{I} \times \mathrm{Y}\). It is a morphism of \(\mathrm{X} \times \mathbf{I}\)-spaces, and even an isomorphism if \(\varphi\) is. If \((x,t,y) \in \mathrm{A} \times \mathbf{I} \times \mathrm{Y}\), we have \(r(x,t) = (x,t)\), whence

\[
h (x, t, y) = (x, t, \mathrm{pr} _ {2} (\varphi (x, t, y))) = (x, t, g ^ {\prime} (x, t, y)) = g (x, t, y),
\]

which proves that h coincides with g on A × I × Y. Similarly, for \(x \in X\) and \(y \in Y\), we have \(r(x,1) = (x,1)\), therefore

\[
h (x, 1, y) = (x, 1, \operatorname{pr} _ {2} (\varphi (x, 1, y))) = (x, 1, f ^ {\prime} (x, 1, y)) = f (x, 1, y)
\]

so that h coincides with f on \(X \times \{1\} \times Y\). The corollary is thus proved.

